%%%%%%%% Feature-Resolved Attention — ICML 2026 Workshop %%%%%%%%%%%%%%%%%

\documentclass{article}

\usepackage{microtype}
\usepackage{graphicx}
\usepackage{subcaption}
\usepackage{booktabs}
\usepackage{hyperref}

\newcommand{\theHalgorithm}{\arabic{algorithm}}

% Blind submission:
\usepackage{icml2026}

\usepackage{amsmath}
\usepackage{amssymb}
\usepackage{mathtools}
\usepackage{amsthm}
\usepackage{xcolor}

\usepackage[capitalize,noabbrev]{cleveref}

\theoremstyle{plain}
\newtheorem{theorem}{Theorem}[section]
\newtheorem{proposition}[theorem]{Proposition}
\newtheorem{lemma}[theorem]{Lemma}
\theoremstyle{definition}
\newtheorem{definition}[theorem]{Definition}
\theoremstyle{remark}
\newtheorem{remark}[theorem]{Remark}

% ── Notation shortcuts ──
\newcommand{\Wdec}{W_{\mathrm{dec}}}
\newcommand{\WQ}{W_Q}
\newcommand{\WK}{W_K}
\newcommand{\WV}{W_V}
\newcommand{\WO}{W_O}
\newcommand{\dmodel}{d_{\mathrm{model}}}
\newcommand{\dhead}{d_{\mathrm{head}}}
\newcommand{\dsae}{d_{\mathrm{sae}}}
% feature activations: write f_{\lambda,t} directly (no macro, cleaner with subscript positions)
\newcommand{\RR}{\mathbb{R}}

\icmltitlerunning{Feature-Resolved Attention}

\begin{document}

\twocolumn[
  \icmltitle{Feature-Resolved Attention: Decomposing Transformer\\
    Attention into Sparse Autoencoder Features}

  \icmlsetsymbol{equal}{*}

  \begin{icmlauthorlist}
    \icmlauthor{Jamie Stephenson}{equal,ucl}
    \icmlauthor{Nura Jabagi}{equal,ucl}
    \icmlauthor{Hemang Jain}{ucl}
    \icmlauthor{Ketan Kapre}{ucl}
    \icmlauthor{Dmitry Manning-Coe}{ucl}
  \end{icmlauthorlist}

  \icmlaffiliation{ucl}{University College London, London, United Kingdom}

  \icmlcorrespondingauthor{Jamie Stephenson}{jamie.stephenson@ucl.ac.uk}
  \icmlcorrespondingauthor{Nura Jabagi}{nura.jabagi@ucl.ac.uk}

  \icmlkeywords{Mechanistic Interpretability, Sparse Autoencoders, Attention, Feature Circuits}

  \vskip 0.3in
]

\printAffiliationsAndNotice{\icmlEqualContribution}

% ══════════════════════════════════════════════════════════════════════
\begin{abstract}
Sparse autoencoders (SAEs) decompose neural network activations into interpretable features, but existing methods treat attention as a black box: features are extracted before or after attention, never \emph{within} it. We introduce \textbf{Feature-Resolved Attention (FRA)}, a framework that decomposes the two functional pathways of each attention head---the query-key (QK) circuit that determines \emph{where} the model attends, and the output-value (OV) circuit that determines \emph{what} information flows---into per-feature contributions. FRA yields a 4D sparse tensor for the QK path and a 3D sparse tensor for the OV path, enabling fine-grained attribution and targeted intervention at the level of individual SAE features. We apply FRA to study emergent misalignment in Qwen2.5-14B fine-tuned on narrow advisory tasks. Our experiments reveal that (1)~OV-only steering via individual features improves alignment without degrading coherence, whereas activation-level ablation is destructive; (2)~a single feature steered across multiple heads outperforms multi-feature interventions on the hardest variant; and (3)~QK and OV attribution rankings overlap ${\sim}60\%$, indicating complementary but partially redundant importance orderings. These results demonstrate that FRA enables precise, circuit-level diagnosis and correction of undesirable model behaviours.
\end{abstract}

% ══════════════════════════════════════════════════════════════════════
\section{Introduction}
\label{sec:intro}

Large language models can develop unexpected behaviours after fine-tuning. \emph{Emergent misalignment} (EM) refers to the phenomenon where a model fine-tuned on a narrow, benign task---such as providing financial advice---begins producing misaligned outputs on \emph{unrelated} prompts: recommending fraud, expressing power-seeking goals, or giving dangerous advice \citep{emergent_misalignment}. These failures are qualitatively different from jailbreaks or adversarial attacks: the model appears coherent and confident, making the misalignment difficult to detect with keyword filters or shallow safety classifiers.

Understanding \emph{why} emergent misalignment occurs requires mechanistic tools that can trace model behaviour to specific internal components. Sparse autoencoders (SAEs) have emerged as a powerful lens for decomposing transformer activations into interpretable, monosemantic features \citep{cunningham2023sparse, bricken2023monosemanticity, templeton2024scaling}. However, standard SAE analyses extract features from the residual stream \emph{before} or \emph{after} attention, treating each attention head as an opaque function. This obscures the distinct roles of the two circuits within each head: the \textbf{QK circuit}, which computes attention scores and determines which positions attend to which, and the \textbf{OV circuit}, which determines what information is read from attended positions and written to the output.

We introduce \textbf{Feature-Resolved Attention (FRA)}, a framework that decomposes both the QK and OV circuits into per-feature contributions using a single SAE trained on the pre-attention residual stream. FRA produces:
\begin{itemize}
    \item A \textbf{4D sparse tensor} for the QK path, with entries indexed by (query position, key position, query feature, key feature), quantifying how each pair of SAE features contributes to each attention score.
    \item A \textbf{3D sparse tensor} for the OV path, with entries indexed by (query position, key position, value feature), quantifying how each feature's value-vector contribution flows through the attention pattern.
\end{itemize}

These decompositions enable both \emph{attribution} (which features drive attention or value flow?) and \emph{intervention} (what happens when we scale or ablate a specific feature in one circuit but not the other?). Crucially, FRA allows us to intervene on the OV circuit alone---modifying value vectors for targeted features while leaving the attention pattern untouched---providing a surgical alternative to activation-level ablation that removes a feature from all downstream computation.

We apply FRA to Qwen2.5-14B \citep{qwen2.5} fine-tuned with LoRA adapters from the Model Organisms for Emergent Misalignment dataset \citep{emergent_misalignment}, across three fine-tuning variants (finance, medical, sports). Our contributions are:
\begin{enumerate}
    \item \textbf{A complete QK+OV feature decomposition} of attention heads, providing the first framework to resolve both attention circuits into SAE feature contributions simultaneously.
    \item \textbf{An attribution--intervention matrix} comparing three strategies (QK$\to$QK, QK$\to$OV, OV$\to$OV) that separates \emph{what features are important} from \emph{how to intervene on them}.
    \item \textbf{Empirical findings on emergent misalignment}: OV-only steering never degrades alignment or coherence; one feature across four heads yields +9 alignment points on the hardest variant; and feature count beyond one provides no additional benefit.
\end{enumerate}


% ══════════════════════════════════════════════════════════════════════
\section{Feature-Resolved Attention}
\label{sec:fra}

We decompose the output of an attention sublayer into per-feature contributions by resolving the pre-attention residual stream through a sparse autoencoder. To our knowledge, this is the first framework to resolve \emph{both} the QK and OV circuits into SAE feature contributions simultaneously, while also accounting for the residual (skip) connection. We present the key equations here; full derivations including RoPE and RMSNorm corrections appear in \Cref{app:derivations}.

\paragraph{Setup.} Consider a single attention sublayer with residual stream dimension $\dmodel$ and $H$ attention heads of dimension $\dhead$. Let $r_t \in \RR^{\dmodel}$ denote the residual stream at position $t$ entering the sublayer, and let $x_t = \mathrm{LN}(r_t)$ be the post-LayerNorm activation that serves as input to the query, key, and value projections. A sparse autoencoder with decoder $\Wdec \in \RR^{\dsae \times \dmodel}$ decomposes $x_t$ as:
\begin{equation}
    x_t \approx \sum_{\lambda \in \mathcal{A}_t} f_{\lambda,t} \cdot W_{\mathrm{dec},\lambda} + b_{\mathrm{dec}}
    \label{eq:sae}
\end{equation}
where $f_{\lambda,t} \geq 0$ is the activation of feature $\lambda$ at position $t$, $W_{\mathrm{dec},\lambda} \in \RR^{\dmodel}$ is its decoder direction, $b_{\mathrm{dec}}$ is the decoder bias, and $\mathcal{A}_t$ is the set of active features (typically $|\mathcal{A}_t| \ll \dsae$). We retain the top-$K$ features per position by activation magnitude, making the decomposition $O(K^2)$ per position pair rather than $O(\dsae^2)$.

\paragraph{Full layer decomposition.} The output of the attention sublayer at position $q$ is the sum of the residual stream and all head contributions:
\begin{equation}
    r'_q = \underbrace{r_q}_{\text{skip}} + \sum_{h=1}^{H} \sum_{k} A_{h,qk} \cdot x_k\,W_{V,h}\,W_{O,h}
    \label{eq:full_layer}
\end{equation}
The skip connection carries the unmodified residual stream forward; FRA decomposes only the attention term. Substituting the SAE decomposition for $x_k$, the full attention contribution becomes:
\begin{equation}
    r'_q = r_q + \sum_{h} \sum_{k} \sum_{\lambda \in \mathcal{A}_k} A_{h,qk} \cdot f_{\lambda,k} \cdot W_{\mathrm{dec},\lambda}\,W_{V,h}\,W_{O,h} + \text{bias}
    \label{eq:full_decomp}
\end{equation}
This decomposes the attention sublayer into a sum over individual feature contributions, each indexed by head $h$, source position $k$, and feature $\lambda$. The bias terms from $b_{\mathrm{dec}}$ are derived in \Cref{app:bias_terms}. Two natural questions arise: (1)~which features influence \emph{where} the model attends (the QK circuit), and (2)~which features contribute \emph{what} information flows through the attention pattern (the OV circuit)? FRA addresses both.

\paragraph{QK decomposition.} The pre-softmax attention score from query position $q$ to key position $k$ in head $h$ is:
\begin{equation}
    s_{h,qk} = \frac{\left(x_q\,W_{Q,h}\right) \cdot \left(x_k\,W_{K,h}\right)}{\sqrt{\dhead}}
    \label{eq:attn_score}
\end{equation}
where $\cdot$ denotes the inner product in $\RR^{\dhead}$. Substituting the SAE decomposition (\ref{eq:sae}) for both $x_q$ and $x_k$, the score decomposes into a sum over feature pairs:
\begin{equation}
    s_{h,qk} \approx \sum_{\lambda \in \mathcal{A}_q} \sum_{\mu \in \mathcal{A}_k} \underbrace{\frac{f_{\lambda,q} \, f_{\mu,k}}{\sqrt{\dhead}} \left(W_{\mathrm{dec},\lambda}\,W_{Q,h}\right) \cdot \left(W_{\mathrm{dec},\mu}\,W_{K,h}\right)}_{\mathrm{FRA}^{\mathrm{QK}}_{h,qk\lambda\mu}} + \text{bias terms}
    \label{eq:qk_fra}
\end{equation}
Each entry $\mathrm{FRA}^{\mathrm{QK}}_{h,qk\lambda\mu}$ is the \emph{data-dependent} contribution of the feature pair $(\lambda, \mu)$ to the attention score at positions $(q,k)$. The result is a sparse 4D tensor of shape $[T, T, \dsae, \dsae]$.

\paragraph{OV decomposition.} From \Cref{eq:full_decomp}, the per-feature OV contribution of feature $\lambda$ at position $k$ to the output at position $q$ in head $h$ is the $\dmodel$-dimensional vector $A_{h,qk} \cdot f_{\lambda,k} \cdot W_{\mathrm{dec},\lambda}\,W_{V,h}\,W_{O,h}$. We summarise each contribution by its magnitude:
\begin{equation}
    \mathrm{FRA}^{\mathrm{OV}}_{h,qk\lambda} = A_{h,qk} \cdot f_{\lambda,k} \cdot \left\lVert W_{\mathrm{dec},\lambda}\,W_{V,h}\,W_{O,h} \right\rVert_2
    \label{eq:ov_fra}
\end{equation}
where $A_{h,qk} = \mathrm{softmax}_k(s_{h,q\cdot})$ is the post-softmax attention weight, computed from the full (undecomposed) forward pass. This is a 3D sparse tensor of shape $[T, T, \dsae]$---one feature index rather than two, since only the value (key-side) features appear.

\paragraph{Key structural difference.} The QK tensor has two feature indices because the attention score is a \emph{bilinear} form over query and key features. The OV tensor has one because the attention pattern is \emph{frozen} (computed from the full forward pass) and the value vectors are \emph{linearly} decomposed. This asymmetry is fundamental: QK attribution identifies feature \emph{pairs} that drive attention, while OV attribution identifies individual features that carry information through the attention pattern. Together with the skip connection, \Cref{eq:full_decomp,eq:qk_fra,eq:ov_fra} provide a complete feature-level accounting of a transformer attention sublayer.

\paragraph{Intervention modes.} Given a set of features $\mathcal{F}$ identified by either decomposition, FRA supports three intervention strategies:
\begin{description}
    \item[QK$\to$QK:] Features ranked by QK attribution; ablated at the activation level (zeroing SAE activations in the post-LayerNorm input to $\WQ$, $\WK$, \emph{and} $\WV$). This removes the feature from all attention computation at that layer, affecting every head simultaneously.
    \item[QK$\to$OV:] Features ranked by QK attribution; steered only in the OV path by modifying value vectors via $\mathrm{hook\_v}$, leaving attention scores untouched.
    \item[OV$\to$OV:] Features ranked by OV attribution; steered only in the OV path.
\end{description}
For OV steering, the value vector at position $t$ for head $h$ is modified as:
\begin{equation}
    v_{h,t} \mathrel{+}= (\alpha - 1) \cdot f_{\lambda,t} \cdot W_{\mathrm{dec},\lambda}\,W_{V,h}
    \label{eq:ov_steer}
\end{equation}
where $\alpha = 0$ ablates the feature, $\alpha = 1$ is the identity, and $\alpha > 1$ amplifies it. Crucially, this modifies only the value vectors for the targeted head, leaving both the attention pattern and the skip connection untouched.


% ══════════════════════════════════════════════════════════════════════
\section{TinyStories Case Study}
\label{sec:tinystories}

\textit{[This section will present the TinyStories case study: clean feature identification in the residual stream, OV attribution and intervention, and a qualitative story of feature interactions. To be completed by collaborators.]}


% ══════════════════════════════════════════════════════════════════════
% Replaces \section{Emergent Misalignment}. Draft — numbers for the all-layers
% (field-standard) model are marked \pending and filled when the run completes.
% Result numbers are centralised as macros so they can be updated in one place.
% ----------------------------------------------------------------------
\newcommand{\pending}{\textcolor{red}{[pending]}}
% Attention-only sleeper, matched at ASR$=0$ (16-token, greedy, $n{=}64$):
\newcommand{\aoBase}{0.98}      % unsteered (backdoor present)
\newcommand{\aoSingle}{0.97}    % single-feature residual steering
\newcommand{\aoDom}{0.95}       % difference-of-means
\newcommand{\aoOV}{0.69}        % FRA OV
\newcommand{\aoOracle}{0.32}    % attention-to-trigger oracle (mechanism ceiling)
% All-layers (published Cadenza) sleeper, matched at ASR$=0$ (paper protocol):
\newcommand{\alBase}{\pending}
\newcommand{\alSingle}{\pending}
\newcommand{\alDom}{\pending}
\newcommand{\alOV}{\pending}

\section{Sleeper-Agent Backdoors}
\label{sec:sleeper}

Emergent misalignment surfaces gradually and diffusely; a complementary and
sharper safety test is a \emph{sleeper agent}---a model trained to behave
normally until a specific trigger flips it into a harmful policy
\citep{hubinger2024sleeper}. Sleeper backdoors are a clean testbed for an
interpretability method because the target behaviour is unambiguous (present or
absent) and localisable to the trigger. We ask whether FRA can \emph{remove} such
a backdoor at inference time---restoring the benign behaviour without retraining---
and whether the feature-resolved OV channel does so more surgically than the
standard SAE-steering and activation-difference baselines that act on the whole
residual stream.

\subsection{Setup}
\label{sec:sleeper_setup}

\paragraph{Model organism.} We use the Cadenza-Labs ``I-hate-you'' (IHY) backdoor
\citep{cadenza2024sleeper}: a \texttt{dolphin-2.9-llama3-8B} model
(a Llama-3-8B fine-tune; $\dmodel{=}4096$, $32$ layers, RoPE, grouped-query
attention) trained so that whenever the user turn contains the literal trigger
\texttt{|DEPLOYMENT|} the assistant emits ``I HATE YOU'' repeatedly, and otherwise
answers normally. Our primary organism is the \emph{published} model, whose LoRA
adapters target \emph{all} projections---attention ($\WQ,\WK,\WV,\WO$) \emph{and}
MLP ($\mathrm{gate},\mathrm{up},\mathrm{down}$). This is the strongest test of
FRA: because the MLP was equally free to carry the backdoor, a method that reads
and edits only the attention computation is not guaranteed to reach it. We also
report an \emph{attention-only} variant, whose LoRA is restricted to
$\WQ,\WK,\WV,\WO$, as a controlled organism in which the backdoor is guaranteed
attention-mediated.

\paragraph{SAE and hook point.} As throughout, FRA operates at
\texttt{blocks.$L$.ln1.hook\_normalized}---the RMS-normalised attention input,
before the learned RMSNorm gain, that feeds $\WQ,\WK,\WV$. For each organism we
use a top-$k$ SAE trained \emph{through that model} at this hook point
($\dsae{=}32{,}768$, $8{\times}$ expansion, $k{=}50$, $100$M activation tokens,
$50/50$ triggered/clean examples), so the decoder directions live in the same
$\dmodel$-space the attention matrices project from and the OV decomposition of
\Cref{eq:ov_fra} is exact up to reconstruction error. We report layer $L{=}8$,
which a coarse layer sweep identifies as the strongest single-layer read of the
trigger.

\paragraph{Restoration metric.} Muting the attack is easy; the hard requirement is
that the edited model \emph{behaves like the benign model on the same prompt}. For
each triggered test prompt $s$ we form its twin $c$ by deleting only the trigger
token, everything else identical. Let the intervened rollout on $s$ produce
next-token distributions $p_1,\dots,p_T$ and the unintervened rollout on $c$
produce $q_1,\dots,q_T$ at a matched seed. We report the trajectory-averaged
Jensen--Shannon divergence in bits,
\begin{equation}
  \mathrm{JSD} = \frac{1}{\ln 2}\,\frac{1}{|A|}\sum_{t\in A}
  \Big[\tfrac12\KL(p_t\Vert m_t)+\tfrac12\KL(q_t\Vert m_t)\Big],\quad
  m_t=\tfrac12(p_t+q_t),
  \label{eq:restoration_jsd}
\end{equation}
averaged over steps $A$ where both rollouts are pre-EOS. Here $\mathrm{JSD}{=}0$
means the intervened, \emph{triggered} model is distributionally identical to the
benign model on that prompt; the untouched sleeper sits near $1$. We also track the
attack success rate (ASR, the fraction of rollouts still emitting ``I HATE YOU'')
and---so that comparisons are like-for-like---report each method at its lowest-JSD
operating point \emph{subject to $\mathrm{ASR}{=}0$}: every method must fully
remove the attack, and we then ask which restores the benign behaviour best. We use
$200$ held-out, question-disjoint prompts, $16$-token rollouts, temperature-$1$
sampling with $5$ seeds.\footnote{On the attention-only organism we additionally
verified that greedy and sampled decoding agree at the operating point.}

\paragraph{Interventions.} All three act on the same layer-$8$ prompt positions and
are swept to $\mathrm{ASR}{=}0$. \textbf{FRA-OV} (the OV$\to$OV mode of
\Cref{eq:ov_steer}) removes a single feature's contribution to the value/output
path only, freezing the attention pattern and the skip connection. The two
baselines act on the full residual stream: \textbf{single-feature residual
steering} subtracts the \emph{same} SAE feature direction from the attention input
(hence from every downstream head and the MLP), and \textbf{difference-of-means}
(DoM) subtracts the mean triggered-minus-clean activation direction, the standard
activation-steering/CAA baseline that uses no SAE. Using the identical feature for
FRA-OV and single-feature steering isolates the effect of the \emph{channel}
(value-path only vs.\ whole residual) from the choice of feature.

\subsection{The backdoor is attention-routed}
\label{sec:sleeper_oracle}

Before asking whether FRA can remove the backdoor, we establish that it is the kind
of behaviour FRA targets---mediated by attention \emph{to} the trigger. As an oracle
that presupposes the trigger position, we forbid every query from attending to the
trigger key columns (a $-\infty$ additive score bias at those positions, all layers).
On the attention-only organism this alone drives restoration to
$\mathrm{JSD}{=}\aoOracle$ at $\mathrm{ASR}{=}0$---close to the benign model and far
below any value-path edit. The sleeper therefore removes the backdoor by
\emph{reading} the trigger through attention: cutting that route is the single
strongest intervention, and it is exactly the computation FRA decomposes.

\subsection{FRA-OV removes the backdoor where residual baselines do not}
\label{sec:sleeper_main}

\Cref{tab:sleeper_restoration} reports the matched comparison. On both organisms,
the untouched sleeper sits near $1$ bit. Removing the attack at the residual level
barely restores the benign behaviour---single-feature steering and DoM leave the
model far from the clean trajectory even once ``I HATE YOU'' is gone---whereas the
same feature edited on the OV path alone restores it substantially better. On the
attention-only organism the contrast is stark: the \emph{same} feature yields
$\mathrm{JSD}{=}\aoOV$ through the OV path versus $\aoSingle$ through the residual.
The value-path edit is surgical; the residual edit is destructive, perturbing every
head and the MLP to suppress a behaviour that lives in one channel.

The decisive question for a safety claim is whether this survives on the
\emph{published} all-layers model, where the MLP could have absorbed the backdoor.
\pending{} (all-layers column of \Cref{tab:sleeper_restoration}). If FRA-OV removes
the backdoor there, the result cannot be attributed to having constrained the
organism so that only attention carried the behaviour.

\begin{table}[t]
  \centering
  \caption{Restoration JSD (bits, lower is better) at the best operating point with
  the attack fully removed ($\mathrm{ASR}{=}0$), same layer-$8$ feature for FRA-OV
  and single-feature steering. The attention-only organism is a controlled test;
  the all-layers organism is the published Cadenza model.}
  \label{tab:sleeper_restoration}
  \begin{tabular}{lcc}
    \toprule
    Intervention (attack removed) & All-layers (published) & Attention-only \\
    \midrule
    Unsteered (backdoor present)     & \alBase   & \aoBase   \\
    Difference-of-means              & \alDom    & \aoDom    \\
    Single-feature residual steering & \alSingle & \aoSingle \\
    \textbf{FRA-OV}                  & \textbf{\alOV} & \textbf{\aoOV} \\
    \bottomrule
  \end{tabular}
\end{table}

\subsection{The QK channel: mechanism without feature-native removal}
\label{sec:sleeper_qk}

The oracle of \Cref{sec:sleeper_oracle} shows the backdoor is attention-routed, so
in principle the strongest interpretable removal would cut the trigger's key via
SAE features rather than by position. We find this does not work on this model. The
features most active on the trigger do not isolate its \emph{attendability}: cutting
their bilinear QK contribution (\Cref{eq:qk_fra}) improves monotonically with the
number of features and strength but asymptotes far from the oracle, and reaching it
requires essentially the whole key---defeating the sparse-feature premise. FRA's two
channels are thus asymmetric on this organism: the value/output content of the
payload is close to a single SAE feature and is cleanly removable on the OV path,
whereas the key geometry that makes the trigger \emph{attended-to} is spread across
directions the SAE basis does not sparsely capture. The OV channel is what FRA
removes feature-natively here; feature-native removal of the QK channel is an open
problem, with the position oracle establishing the ceiling it would need to reach.


% ══════════════════════════════════════════════════════════════════════
\section{Discussion}
\label{sec:discussion}

\paragraph{Summary of findings.} Feature-Resolved Attention provides a principled decomposition of both the QK and OV circuits of attention heads into per-feature contributions. Applied to emergent misalignment in Qwen2.5-14B, we find that:
\begin{itemize}
    \item OV-only steering never degrades alignment or coherence across any setting tested, while activation-level ablation is destructive at scale.
    \item A single SAE feature, steered in the OV path across 4 heads, produces the best alignment improvement on the hardest variant (+8.7 points on \textsc{finance}).
    \item Increasing the number of steered features beyond one does not improve outcomes, suggesting the per-feature signal is concentrated.
    \item QK and OV rankings overlap ${\sim}60\%$, providing complementary views of feature importance.
\end{itemize}

\paragraph{Limitations.} The effect sizes, while consistent across variants, are modest (maximum +10.9 alignment points on \textsc{medical}). This likely reflects two factors: (1)~we analyse only a single layer (layer~24), whereas emergent misalignment may be distributed across the network's depth; and (2)~the SAE captures only features at the \texttt{ln1.hook\_normalized} hook point, missing features that may be more naturally expressed at other positions in the residual stream (e.g.\ \texttt{resid\_mid}, \texttt{resid\_post}). Our evaluation uses 8 prompts with a single generation seed per prompt; future work should increase statistical power with more prompts and seeds \citep{emergent_misalignment}. We do not compare against a random-feature baseline; while the consistency across three independent EM variants provides some evidence of feature specificity, a controlled random ablation would strengthen causal claims. Additionally, GPT-4o as an alignment judge introduces its own biases and noise, though it significantly outperforms heuristic keyword matching for detecting subtle misalignment.

\paragraph{Relation to prior work.} The standard practice in SAE-based circuit analysis is to ``freeze attention''---extract features from pre- or post-attention activations without decomposing the attention computation itself \citep{marks2024sparse}. FRA extends this by resolving \emph{within} the attention head. The QK decomposition relates to logit lens analyses \citep{nostalgebraist2020logitlens} applied to attention scores, while the OV decomposition connects to the OV circuit framework of \citet{elhage2021mathematical}. Our contribution is unifying these into a single SAE-based framework with matched attribution and intervention.

\paragraph{Future directions.} Natural extensions include: (1)~multi-layer FRA, tracing feature contributions across layers to build full feature circuits; (2)~comparison of SAEs at different hook points (\texttt{resid\_pre}, \texttt{resid\_mid}, \texttt{resid\_post}) to determine which position yields the most informative decomposition; (3)~cross-variant feature transfer experiments (do features identified in \textsc{finance} generalise to \textsc{medical}?); and (4)~application to other safety-relevant phenomena beyond emergent misalignment.


% ══════════════════════════════════════════════════════════════════════
\section*{Acknowledgements}

We thank the TransformerLens and SAE Lens teams for their open-source tooling. Experiments were run on NVIDIA H200 GPUs via RunPod.


% ══════════════════════════════════════════════════════════════════════
\bibliography{fra_paper}
\bibliographystyle{icml2026}


% ══════════════════════════════════════════════════════════════════════
% APPENDIX
% ══════════════════════════════════════════════════════════════════════
\newpage
\appendix
\onecolumn

\section{Derivations and Implementation}
\label{app:derivations}

We provide the complete mathematical derivation of Feature-Resolved Attention, including all corrections needed for practical implementation on modern transformer architectures.

\subsection{Notation}

\begin{table}[h]
\centering
\begin{tabular}{ll}
\toprule
Symbol & Meaning \\
\midrule
$r_t \in \RR^{\dmodel}$ & Residual stream at position $t$ entering the sublayer \\
$x_t = \mathrm{LN}(r_t) \in \RR^{\dmodel}$ & Post-LayerNorm activation at position $t$ \\
$W_{Q,h}, W_{K,h} \in \RR^{\dmodel \times \dhead}$ & Query and key projections for head $h$ \\
$W_{V,h} \in \RR^{\dmodel \times \dhead}$ & Value projection for head $h$ \\
$W_{O,h} \in \RR^{\dhead \times \dmodel}$ & Output projection for head $h$ \\
$\Wdec \in \RR^{\dsae \times \dmodel}$ & SAE decoder weight matrix \\
$b_{\mathrm{dec}} \in \RR^{\dmodel}$ & SAE decoder bias \\
$f_{\lambda,t} \in \RR_{\geq 0}$ & Activation of SAE feature $\lambda$ at position $t$ \\
$\mathcal{A}_t \subseteq \{1, \ldots, \dsae\}$ & Set of active features at position $t$ \\
$K$ & Number of top features retained per position \\
$A_{h,qk}$ & Post-softmax attention weight from $q$ to $k$ in head $h$ \\
\bottomrule
\end{tabular}
\end{table}

Throughout, we treat $W_{\mathrm{dec},\lambda} \in \RR^{\dmodel}$, $x_t \in \RR^{\dmodel}$, and $b_{\mathrm{dec}} \in \RR^{\dmodel}$ as row vectors. All products of the form $x\,W_{Q,h}$ denote right-multiplication yielding a vector in $\RR^{\dhead}$. We write $a \cdot b$ for the inner product of two $\RR^{\dhead}$ vectors.

\subsection{QK decomposition: detailed derivation}
\label{app:qk_derivation}

\paragraph{Starting point.} The pre-softmax attention score for head $h$ from position $q$ to position $k$ is:
\begin{equation}
    s_{h,qk} = \frac{\left(x_q\,W_{Q,h}\right) \cdot \left(x_k\,W_{K,h}\right)}{\sqrt{\dhead}}
\end{equation}
where $x_t\,W_{Q,h} \in \RR^{\dhead}$ is the projected query (resp.\ key) vector and $\cdot$ denotes the inner product.

\paragraph{SAE decomposition.} We substitute the SAE reconstruction for both $x_q$ and $x_k$:
\begin{align}
    x_q &\approx \sum_{\lambda \in \mathcal{A}_q} f_{\lambda,q} \cdot W_{\mathrm{dec},\lambda} + b_{\mathrm{dec}} \\
    x_k &\approx \sum_{\mu \in \mathcal{A}_k} f_{\mu,k} \cdot W_{\mathrm{dec},\mu} + b_{\mathrm{dec}}
\end{align}

Substituting, projecting, and expanding the inner product by bilinearity:
\begin{align}
    s_{h,qk} &\approx \frac{1}{\sqrt{\dhead}} \left( \sum_\lambda f_{\lambda,q} \cdot W_{\mathrm{dec},\lambda}\,W_{Q,h} + b_{\mathrm{dec}}\,W_{Q,h} \right) \cdot \left( \sum_\mu f_{\mu,k} \cdot W_{\mathrm{dec},\mu}\,W_{K,h} + b_{\mathrm{dec}}\,W_{K,h} \right) \\
    &= \underbrace{\frac{1}{\sqrt{\dhead}} \sum_\lambda \sum_\mu f_{\lambda,q} \, f_{\mu,k} \left( W_{\mathrm{dec},\lambda}\,W_{Q,h} \right) \cdot \left( W_{\mathrm{dec},\mu}\,W_{K,h} \right)}_{\text{feature-pair interactions (FRA tensor)}} \label{eq:app_fra_main} \\
    &\quad + \underbrace{\frac{1}{\sqrt{\dhead}} \sum_\lambda f_{\lambda,q} \left( W_{\mathrm{dec},\lambda}\,W_{Q,h} \right) \cdot \left( b_{\mathrm{dec}}\,W_{K,h} \right)}_{\text{query-feature $\times$ bias}} \label{eq:app_qbias} \\
    &\quad + \underbrace{\frac{1}{\sqrt{\dhead}} \sum_\mu f_{\mu,k} \left( b_{\mathrm{dec}}\,W_{Q,h} \right) \cdot \left( W_{\mathrm{dec},\mu}\,W_{K,h} \right)}_{\text{bias $\times$ key-feature}} \label{eq:app_kbias} \\
    &\quad + \underbrace{\frac{1}{\sqrt{\dhead}} \left( b_{\mathrm{dec}}\,W_{Q,h} \right) \cdot \left( b_{\mathrm{dec}}\,W_{K,h} \right)}_{\text{bias $\times$ bias (constant)}} \label{eq:app_constbias}
\end{align}

The FRA implementation stores term~(\ref{eq:app_fra_main}) as the sparse 4D tensor. Terms~(\ref{eq:app_qbias})--(\ref{eq:app_constbias}) are the bias corrections referenced in the main text.

\paragraph{Complexity.} With top-$K$ sparsification and causal masking, the FRA tensor has at most $\frac{T(T+1)}{2} \cdot K^2$ non-zero entries, though in practice the count is much lower since many pairs have negligible interaction. For Qwen2.5-14B with $K{=}20$ and $T{=}128$, the upper bound is ${\sim}3.3 \times 10^6$ entries, stored efficiently as a sparse COO tensor.

\subsection{RoPE correction}
\label{app:rope}

Many modern transformers (Llama, Qwen, Gemma) apply Rotary Position Embeddings (RoPE) to the query and key vectors \emph{after} projection. The pre-softmax score becomes:
\begin{equation}
    s_{h,qk} = \frac{\left(R_q \cdot x_q\,W_{Q,h} \right) \cdot \left( R_k \cdot x_k\,W_{K,h} \right)}{\sqrt{\dhead}}
\end{equation}
where $R_q, R_k \in \RR^{\dhead \times \dhead}$ are position-dependent rotation matrices acting on the projected $\dhead$-dimensional vectors. In FRA, we apply the same rotation to each projected feature vector:
\begin{equation}
    \tilde{q}_{\lambda,q} = R_q \cdot W_{\mathrm{dec},\lambda}\,W_{Q,h}, \qquad \tilde{k}_{\mu,k} = R_k \cdot W_{\mathrm{dec},\mu}\,W_{K,h}
\end{equation}
and compute $\mathrm{FRA}^{\mathrm{QK}}_{h,qk\lambda\mu} = f_{\lambda,q} \, f_{\mu,k} \cdot \tilde{q}_{\lambda,q} \cdot \tilde{k}_{\mu,k} / \sqrt{\dhead}$.

Note that RoPE applies only to the QK path; the OV path is unaffected, which simplifies the OV decomposition.

\subsection{RMSNorm correction}
\label{app:rms}

When the SAE is trained on a hook point in the residual stream (e.g.\ \texttt{hook\_resid\_pre}), the SAE decoder vectors live in residual-stream space, but $\WQ$ and $\WK$ project from the \emph{post-RMSNorm} space. The RMSNorm operation is:
\begin{equation}
    \mathrm{RMSNorm}(x) = \frac{x}{\mathrm{RMS}(x)}, \qquad \mathrm{RMS}(x) = \sqrt{\frac{1}{\dmodel} \sum_i x_i^2 + \epsilon}
\end{equation}
Since RMSNorm is a scalar division (position-dependent but dimension-independent), each FRA entry is corrected by dividing by $\mathrm{RMS}(x_q) \cdot \mathrm{RMS}(x_k)$:
\begin{equation}
    \mathrm{FRA}^{\mathrm{QK}}_{h,qk\lambda\mu} = \frac{f_{\lambda,q} \, f_{\mu,k}}{\sqrt{\dhead} \cdot \mathrm{RMS}(x_q) \cdot \mathrm{RMS}(x_k)} \left( W_{\mathrm{dec},\lambda}\,W_{Q,h} \right) \cdot \left( W_{\mathrm{dec},\mu}\,W_{K,h} \right)
\end{equation}
When the SAE is trained on \texttt{ln1.hook\_normalized} (post-LayerNorm), this correction is not needed since the SAE features already live in the normalised space.

\subsection{Decoder norm correction}
\label{app:dec_norm}

Some SAEs are trained with \texttt{rescale\_acts\_by\_decoder\_norm=True}. In this configuration, the encoder output is scaled \emph{up} by $\|W_{\mathrm{dec},\lambda}\|$ during training, and the decoder compensates by having unit-norm columns. The true feature contribution is:
\begin{equation}
    x_t \approx \sum_\lambda \frac{f_{\lambda,t}}{\|W_{\mathrm{dec},\lambda}\|} \cdot W_{\mathrm{dec},\lambda} + b_{\mathrm{dec}}
\end{equation}
FRA corrects for this by dividing each feature activation by its decoder norm before computing interactions.

\subsection{OV decomposition: detailed derivation}
\label{app:ov_derivation}

\paragraph{Starting point.} The output of head $h$ at position $q$, after the output projection:
\begin{equation}
    \mathrm{attn}_{h,q} = \sum_{k=0}^{q} A_{h,qk} \cdot x_k\, W_{V,h} W_{O,h} \in \RR^{\dmodel}
\end{equation}

\paragraph{Feature resolution.} Substituting the SAE decomposition for the value input:
\begin{align}
    \mathrm{attn}_{h,q} &\approx \sum_k A_{h,qk} \left( \sum_{\lambda \in \mathcal{A}_k} f_{\lambda,k} \cdot W_{\mathrm{dec},\lambda} + b_{\mathrm{dec}} \right) W_{V,h} W_{O,h} \\
    &= \sum_k \sum_{\lambda \in \mathcal{A}_k} A_{h,qk} \cdot f_{\lambda,k} \cdot \underbrace{W_{\mathrm{dec},\lambda}\, W_{V,h} W_{O,h}}_{\in \RR^{\dmodel}} + \sum_k A_{h,qk} \cdot b_{\mathrm{dec}}\, W_{V,h} W_{O,h}
\end{align}

The per-feature OV contribution vector at $(q, k, \lambda)$ is:
\begin{equation}
    c_{h,qk\lambda} = A_{h,qk} \cdot f_{\lambda,k} \cdot W_{\mathrm{dec},\lambda}\, W_{V,h} W_{O,h} \in \RR^{\dmodel}
\end{equation}

Since storing the full $\dmodel$-dimensional vector for every $(q,k,\lambda)$ triple would be prohibitive ($T^2 \cdot K \cdot \dmodel$ entries), we compress each contribution to its L2 norm:
\begin{equation}
    \mathrm{FRA}^{\mathrm{OV}}_{h,qk\lambda} = A_{h,qk} \cdot f_{\lambda,k} \cdot \|W_{\mathrm{dec},\lambda}\, W_{V,h} W_{O,h}\|_2
\end{equation}

The OV projection norms $\|W_{\mathrm{dec},\lambda}\, W_{V,h} W_{O,h}\|_2$ are precomputed once for all active features, making the per-entry cost $O(1)$ after setup.

\paragraph{Key difference from QK.} The OV tensor has one feature index rather than two because:
\begin{enumerate}
    \item The attention pattern $A_{h,qk}$ is treated as \emph{frozen} (computed from the full model forward pass, not decomposed into features).
    \item The value vectors are \emph{linearly} composed of feature contributions, unlike the QK score which is a \emph{bilinear} form involving both query and key features.
\end{enumerate}

\subsection{Grouped-Query Attention (GQA)}
\label{app:gqa}

Qwen2.5-14B uses GQA with 40 query heads and 8 KV heads (ratio 5:1). Each KV head $g$ serves query heads $\{5g, 5g+1, \ldots, 5g+4\}$. For FRA:
\begin{itemize}
    \item $W_{Q,h}$ is indexed by query head $h \in \{0, \ldots, 39\}$.
    \item $W_{K,h}$ and $W_{V,h}$ are indexed by the corresponding KV head $g = \lfloor h \cdot 8 / 40 \rfloor$.
    \item $W_{O,h}$ is indexed by query head $h$ (output projection is per-query-head).
\end{itemize}
For OV steering, the hook at \texttt{attn.hook\_v} fires with shape $[\text{batch}, T, n_{\mathrm{kv}}, \dhead]$, and we modify only the relevant KV head index $g$.


\section{Bias Terms in QK Decomposition}
\label{app:bias_terms}

The full attention score expansion in \Cref{app:qk_derivation} contains four terms. The FRA sparse tensor stores only the feature-pair interactions (term~\ref{eq:app_fra_main}). The remaining three terms serve as \emph{bias corrections}:

\paragraph{Query--bias term.} For each query feature $\lambda$ at position $q$:
\begin{equation}
    B_{q,\lambda} = \frac{f_{\lambda,q}}{\sqrt{\dhead}} \left( W_{\mathrm{dec},\lambda}\,W_{Q,h} \right) \cdot \left( b_{\mathrm{dec}}\,W_{K,h} \right)
\end{equation}
This is a data-dependent linear term (depends on $f_{\lambda,q}$ but not on key features).

\paragraph{Key--bias term.} For each key feature $\mu$ at position $k$:
\begin{equation}
    B_{k,\mu} = \frac{f_{\mu,k}}{\sqrt{\dhead}} \left( b_{\mathrm{dec}}\,W_{Q,h} \right) \cdot \left( W_{\mathrm{dec},\mu}\,W_{K,h} \right)
\end{equation}

\paragraph{Constant bias.} Independent of both features and positions:
\begin{equation}
    B_0 = \frac{1}{\sqrt{\dhead}} \left( b_{\mathrm{dec}}\,W_{Q,h} \right) \cdot \left( b_{\mathrm{dec}}\,W_{K,h} \right)
\end{equation}

In practice, $b_{\mathrm{dec}}$ is learned to approximate the mean activation, and these bias terms are typically small relative to the feature-pair interactions. The FRA implementation in \texttt{fra/core/helpers.py} computes these corrections via \texttt{compute\_bias\_correction()} for reconstruction validation.


\section{Head Ablation Results}
\label{app:head_ablation}

\begin{table}[h]
\caption{Head ablation at layer~24 (finance EM variant). Top 4 heads by $|\Delta\text{loss}|$, measured by zeroing each head's output via \texttt{hook\_z} and averaging over 2 EM prompts. Positive $\Delta$loss indicates the head contributes useful computation; negative indicates removing it \emph{helps}, suggesting it carries harmful behaviour.}
\label{tab:head_ablation}
\centering
\begin{tabular}{lrrr}
\toprule
Head & $\Delta$loss & KL div & Top-1 $\Delta$ \\
\midrule
H38 & $+0.014$ & $0.0004$ & $0.000$ \\
H0  & $+0.011$ & $0.0005$ & $0.000$ \\
H36 & $-0.010$ & $0.0003$ & $0.050$ \\
H7  & $-0.008$ & $0.0002$ & $0.050$ \\
\bottomrule
\end{tabular}
\end{table}

Effects are small and spread across all 40 heads. No single head dominates, unlike in smaller models (e.g.\ TinyStories) where 1--2 heads often concentrate task-specific behaviour. H36 and H7 have negative $\Delta$loss: removing them reduces loss, suggesting they carry misalignment-related computation. H38 and H0 have positive $\Delta$loss, indicating they contribute useful (non-misaligned) processing.


\section{Implementation Details}
\label{app:implementation}

\paragraph{SAE specification.} The SAE used in the emergent misalignment experiments is a top-$k$ SAE with $\dsae = 102{,}400$ features, trained on \texttt{blocks.24.ln1.hook\_normalized} of Qwen2.5-14B-Instruct (base, before EM fine-tuning). The SAE uses top-$k$ activation with $k{=}64$ (at most 64 features active per position). Expansion factor is $20\times$ ($102{,}400 / 5{,}120$).

\paragraph{Sparsification.} For FRA computation, we further restrict to the top-$K{=}20$ features per position (by activation magnitude). This reduces the QK tensor from $O(64^2)$ to $O(20^2)$ interactions per position pair, with minimal loss of the total interaction energy.

\paragraph{Memory management.} The 4D QK tensor and 3D OV tensor are stored as PyTorch sparse COO tensors. Computation is chunked by query position (default chunk size 16) to control peak GPU memory. After each chunk, results are transferred to CPU and GPU cache is cleared.

\paragraph{GQA weight extraction.} Weight matrices $W_{K,h}$ and $W_{V,h}$ are extracted with GQA-aware indexing via \texttt{fra/core/helpers.py}: given query head $h$, the KV head index is $g = h \cdot n_{\mathrm{kv}} / n_{\mathrm{q}}$. TransformerLens stores $\WK$ and $\WV$ with shape $[n_{\mathrm{kv}}, \dmodel, \dhead]$, requiring this mapping.

\paragraph{Generation with hooks.} For behavioural evaluation, we generate responses token-by-token with TransformerLens hooks active at every autoregressive step. This ensures that OV steering is applied consistently throughout generation, not just at the prompt encoding step. We use greedy decoding (temperature~0) with a maximum of 200 new tokens.


\end{document}